% 原PDF137页4.B.3标题—144页附录4.C标题前。
\subsection*{4.B.3\quad Cox—Reid调整剖面似然}
\addcontentsline{toc}{subsection}{4.B.3 Cox—Reid调整剖面似然}

常用统计推断以观测似然函数为基础。存在干扰参数时，通过调整剖面似然函数进行统计推断。第4.2A2节假设离散参数已知，但实际中通常未知，需要估计。本节介绍Cox—Reid调整剖面似然，给出离散参数 $\phi$ 的估计方法。为了方便对 $\phi$ 计算对数似然，先改变负二项分布对数似然的形式。

回顾式（4.B.5），
\begin{equation}
\begin{aligned}
P(Y=y_{gi})
&=\frac{\Gamma(y_{gi}+\alpha)}{y_{gi}!\Gamma(\alpha)}
 \left(\frac1{1+\beta}\right)^{y_{gi}}
 \left(\frac\beta{1+\beta}\right)^\alpha\\
&=\frac{\Gamma(y_{gi}+1/\phi_g)}{y_{gi}!\Gamma(1/\phi_g)}
 (\phi_g\mu_{gi})^{y_{gi}}
 \left(\frac1{1+\phi_g\mu_{gi}}\right)^{y_{gi}+1/\phi_g}.
\end{aligned}\tag{4.B.27}
\end{equation}
注意
\begin{equation}
\begin{aligned}
\frac{\Gamma(y_{gi}+1/\phi_g)}{y_{gi}!\Gamma(1/\phi_g)}
&=\frac{\begin{gathered}
 (y_{gi}+1/\phi_g-1)(y_{gi}+1/\phi_g-2)\cdots(1/\phi_g+1)\\
 {}\times(1/\phi_g)(1/\phi_g-1)(1/\phi_g-2)\cdots
 \end{gathered}}{
 y_{gi}!(1/\phi_g-1)(1/\phi_g-2)\cdots}\\
&=\frac{(y_{gi}+1/\phi_g-1)(y_{gi}+1/\phi_g-2)\cdots(1/\phi_g+1)(1/\phi_g)}{y_{gi}!}\\
&=\frac{[(y_{gi}-1)\phi_g+1][(y_{gi}-2)\phi_g+1]\cdots[\phi_g+1]}{
 \phi_g^{y_{gi}}y_{gi}!}.
\end{aligned}\tag{4.B.28}
\end{equation}
将式（4.B.28）代入式（4.B.27），得到
\begin{equation}
\begin{aligned}
P(Y=y_{gi})={}&
\frac{[(y_{gi}-1)\phi_g+1][(y_{gi}-2)\phi_g+1]\cdots[\phi_g+1]}{
 \phi_g^{y_{gi}}y_{gi}!}\\
&{}\times(\phi_g\mu_{gi})^{y_{gi}}
 \left(\frac1{1+\phi_g\mu_{gi}}\right)^{y_{gi}+1/\phi_g}.
\end{aligned}\tag{4.B.29}
\end{equation}
利用式（4.B.29），得到对数似然
\begin{equation}
\begin{aligned}
l(\phi_g;y_g,\widehat\gamma_g)
=\sum_{i=1}^{n}\biggl[&\sum_{l=1}^{y_{gi}-1}\log(1+l\phi_g)
-y_{gi}\log\phi_g-\log y_{gi}!\\
&{}+y_{gi}\log(\phi_g\mu_{gi})
-\left(y_{gi}+\frac1{\phi_g}\right)\log(1+\phi_g\mu_{gi})\biggr].
\end{aligned}\tag{4.B.30}
\end{equation}
$\phi_g$ 的调整剖面似然（adjusted profile likelihood，APL）定义为惩罚对数似然
\begin{equation}
\mathrm{APL}_g(\phi_g,\widehat\gamma_g)
=l(\phi_g;y_g,\widehat\gamma_g)-\frac12\log|I_g|.\tag{4.B.31}
\end{equation}
为计算极大化 $\mathrm{APL}_g(\phi_g)$ 的一阶最优条件，先计算
\[
\frac{\partial l(\phi_g;y_g,\widehat\gamma_g)}{\partial\phi_g}.
\]
由式（4.B.30），
\begin{equation}
\begin{aligned}
\frac{\partial l(\phi_g;y_g,\widehat\gamma_g)}{\partial\phi_g}
&=\sum_{i=1}^{n}\biggl[\sum_{l=1}^{y_{gi}-1}\frac l{1+l\phi_g}
 -\frac{y_{gi}}{\phi_g}+\frac{y_{gi}}{\phi_g}
 +\frac{\log(1+\phi_g\mu_g)}{\phi_g^2}\\
&\hspace{6em}{}-\left(y_{gi}+\frac1{\phi_g}\right)
 \frac{\mu_{gi}}{1+\phi_g\mu_{gi}}\biggr]\\
&=\sum_{i=1}^{n}\left[\sum_{l=1}^{y_{gi}-1}\frac l{1+l\phi_g}
 +\frac{\log(1+\phi_g\mu_g)}{\phi_g^2}
 -\left(y_{gi}+\frac1{\phi_g}\right)\frac{\mu_{gi}}{1+\phi_g\mu_{gi}}\right].
\end{aligned}\tag{4.B.32}
\end{equation}
接着考虑式（4.B.31）的第二项。利用假设 $y_{gi}\approx\mu_{gi}+1/\phi_g$，得到
\begin{equation}
W_1=\frac{\partial W}{\partial\phi_g}=
-\begin{bmatrix}
\dfrac{\mu_{g1}^2}{(1+\phi_g\mu_{g1})^2}&\cdots&0\\
\vdots&\vdots&\vdots\\
0&\cdots&\dfrac{\mu_{gn}^2}{(1+\phi_g\mu_{gn})^2}
\end{bmatrix}.\tag{4.B.33}
\end{equation}
利用Xiong（2018）式（1.160）与链式法则，得到
\begin{equation}
\frac{\partial\log|I_g|}{\partial\phi_g}
=-I_g^{-1}X\frac{\partial W}{\partial\phi_g}X^T
=-I_g^{-1}XW_1X^T.\tag{4.B.34}
\end{equation}
因此，APL关于 $\phi_g$ 的导数为
\begin{equation}
\frac{\partial\mathrm{APL}(\phi_g)}{\partial\phi_g}
=\frac{\partial l(\phi_g;y_g,\widehat\beta_g)}{\partial\phi_g}
+\frac12I_g^{-1}XW_1X^T.\tag{4.B.35}
\end{equation}
下面计算APL关于 $\phi_g$ 的二阶导数。对式（4.B.32）两侧求导，得到
\begin{equation}
\begin{aligned}
\frac{\partial^2l(\phi_g;y_g,\widehat\beta_g)}{\partial\phi_g^2}
={}&\sum_{i=1}^{n}\biggl[-\sum_{l=1}^{y_{gi}-1}\frac{l^2}{(1+l\phi_g)^2}\\
&\hspace{3em}{}+\frac{\mu_g\phi_g-2(1+\phi_g\mu_g)\log(1+\phi_g\mu_g)}{
 (1+\phi_g\mu_g)\phi_g^3}\biggr]\\
&{}+\frac{\mu_{gi}}{(1+\phi_g\mu_g)\phi_{gi}^2}
+(y_{gi}+\phi_g^{-1})\frac{\mu_{gi}^2}{(1+\phi_g\mu_g)^2}.
\end{aligned}\tag{4.B.36}
\end{equation}
利用式（4.B.33），得到
\begin{equation}
W_2=\frac{\partial W_1}{\partial\phi_g}=
\begin{bmatrix}
\dfrac{2\mu_{g1}^3}{(1+\phi_g\mu_{g1})^3}&\cdots&0\\
\vdots&\vdots&\vdots\\
0&\cdots&\dfrac{\mu_{gn}^3}{(1+\phi_g\mu_{gn})^3}
\end{bmatrix}.\tag{4.B.37}
\end{equation}
因此，由式（4.B.34）与式（4.B.35），
\[
\begin{aligned}
\frac{\partial^2\mathrm{APL}(\phi_g)}{\partial\phi_2^2}
={}&\frac{\partial^2l(\phi_g;y_g,\widehat\beta_g)}{\partial\phi_g^2}
-\frac12\frac{\partial I_g^{-1}}{\partial\phi_g}XW_1X^T
+\frac12I_g^{-1}X\frac{\partial W_1}{\partial\phi_g}X^T\\
={}&\frac{\partial^2l(\phi_g;y_g,\widehat\beta_g)}{\partial\phi_g^2}
-\frac12(I_g^{-1}XW_1X^T)^2+\frac12I_g^{-1}XW_2X^T.
\end{aligned}
\]
可以用Newton—Raphson迭代算法估计离散参数 $\phi_g$：
\begin{equation}
\begin{aligned}
\phi_g^{(\mathrm{new})}
&=\phi_g^{(\mathrm{old})}
-\frac{\left.\dfrac{\partial\mathrm{APL}(\phi_g)}{\partial\phi_g}
 \right|_{\phi_g^{(\mathrm{old})}}}{
 \left.\dfrac{\partial^2\mathrm{APL}(\phi_g)}{\partial\phi_2^2}
 \right|_{\phi_g^{(\mathrm{old})}}}\\
&=\phi_g^{(\mathrm{old})}+\delta,
\end{aligned}\tag{4.B.38}
\end{equation}
其中
\[
\delta=-\frac{
\dfrac{\partial l(\phi_g;y_g\widehat\beta_g)}{\partial\phi_g}
+\dfrac12I_g^{-1}XW_1X^T}{
\dfrac{\partial^2l(\phi_g;y_g\widehat\beta_g)}{\partial\phi_g^2}
-\dfrac12(I_g^{-1}XW_1X^T)^2+\dfrac12I_g^{-1}XW_2X^T},
\]
\[
\frac{\partial l(\phi_g;y_g,\widehat\beta_g)}{\partial\phi_g}
=\sum_{i=1}^{n}\left[\sum_{l=1}^{y_{gi}-1}\frac l{1+l\phi_g}
+\frac{\log(1+\phi_g\mu_g)}{\phi_g^2}
-\left(y_{gi}+\frac1{\phi_g}\right)\frac{\mu_{gi}}{1+\phi_g\mu_{gi}}\right],
\]
\[
I_g=XWX^T,\qquad
W=\begin{bmatrix}
\dfrac{\mu_{g1}y_{g1}\phi_g}{(1+\phi_g\mu_{g1})^2}&\cdots&0\\
\vdots&\vdots&\vdots\\
0&\cdots&\dfrac{\mu_{gn}y_{gn}\phi_g}{(1+\phi_g\mu_{gn})^2}
\end{bmatrix},
\]
\[
W_1=\frac{\partial W}{\partial\phi_g}=
-\begin{bmatrix}
\dfrac{\mu_{g1}^2}{(1+\phi_g\mu_{g1})^2}&\cdots&0\\
\vdots&\vdots&\vdots\\
0&\cdots&\dfrac{\mu_{gn}^2}{(1+\phi_g\mu_{gn})^2}
\end{bmatrix},
\]
\[
\begin{aligned}
\frac{\partial^2l(\phi_g;y_g,\widehat\beta_g)}{\partial\phi_g^2}
={}&\sum_{i=1}^{n}\biggl[-\sum_{l=1}^{y_{gi}-1}\frac{l^2}{(1+l\phi_g)^2}\\
&\hspace{3em}{}+\frac{\mu_g\phi_g-2(1+\phi_g\mu_g)\log(1+\phi_g\mu_g)}{
 (1+\phi_g\mu_g)\phi_g^3}\biggr]\\
&{}+\frac{\mu_{gi}}{(1+\phi_g\mu_g)\phi_{gi}^2}
+(y_{gi}+\phi_g^{-1})\frac{\mu_{gi}^2}{(1+\phi_g\mu_g)^2},
\end{aligned}
\]
\[
W_2=\frac{\partial W_1}{\partial\phi_g}=
\begin{bmatrix}
\dfrac{2\mu_{g1}^3}{(1+\phi_g\mu_{g1})^3}&\cdots&0\\
\vdots&\vdots&\vdots\\
0&\cdots&\dfrac{2\mu_{gn}^3}{(1+\phi_g\mu_{gn})^3}
\end{bmatrix},\qquad
\mu_{gi}=N_ie^{x_i^T\widehat\gamma_g}.
\]
Newton—Raphson算法通常使似然增加，但也可能使似然下降。为了保证每次迭代都增加似然，引入松弛常数，将式（4.B.38）化为
\begin{equation}
\phi_g^{(\mathrm{new})}=\phi_g^{(\mathrm{old})}+\alpha\delta,\tag{4.B.39}
\end{equation}
其中 $\alpha$ 为松弛常数。通过线搜索确定 $\alpha$，使
\[
\mathrm{APL}_g(\phi_g^{(\mathrm{new})})\ge\mathrm{APL}_g(\phi_g^{(\mathrm{old})}).
\]
式（4.B.38）、（4.B.39）均假设已经有估计量 $\widehat\gamma_g$，但估计 $\gamma_g$ 又需要离散参数 $\phi_g$，因此需要在二者之间交替迭代。

\subsection*{4.B.4\quad 检验统计量}
\addcontentsline{toc}{subsection}{4.B.4 检验统计量}

我们关心对数线性模型系数 $\gamma_g$ 的显著性。零假设为
\[
H_0:\gamma_g=0.
\]
定义检验统计量（Huang et al. 2015）。不存在过度离散时，检验 $H_0:\gamma_g=0$ 的似然比统计量为
\begin{equation}
\mathrm{LR}_g=2\bigl[l(\widehat\gamma_g,\widehat\phi_g)
-l(\gamma_g^0,\widehat\phi_g)\bigr],\tag{4.B.40}
\end{equation}
其中
\[
\begin{aligned}
l(\widehat\gamma_g,\phi_g)
={}&\left(\sum_{i=1}^{n}y_{gi}\right)\log\phi_g\\
&{}+\sum_{i=1}^{n}\biggl[y_{gi}(x_i^T\widehat\gamma_g+\log N_i)
-\left(y_{gi}+\frac1{\phi_g}\right)
\log(1+\phi_gN_ie^{x_i^T\widehat\gamma_g})\biggr].
\end{aligned}
\]
存在过度离散时，检验相同零假设的似然比统计量定义为
\begin{equation}
\mathrm{LR}_{gd}=\frac{2\bigl[l(\widehat\gamma_g,\widehat\phi_g)
-l(\gamma_g^0,\widehat\phi_g)\bigr]}{\widehat\phi_g}.\tag{4.B.41}
\end{equation}
在零假设下，$\mathrm{LR}_g,\mathrm{LR}_{gd}$ 分别渐近服从中心 $\chi_{(q)}^2$ 与 $F_{(q,n-q-1)}$ 分布。

检验差异表达的一种自然方法，是比较病例与对照的表达差异。考虑A组（病例）与B组（对照）。令 $n_{gA},n_{gB}$ 分别为两组中基因 $g$ 的读段数，$n_A,n_B$ 为两组的读段总数，$n_A-n_{gA},n_B-n_{gB}$ 为其余基因的读段数。
令 $\mu_{gA},\mu_{gB}$ 分别为两组中基因 $g$ 的平均读段数。基因 $g$ 在两组间无表达差异的零假设为
\[
H_0:\mu_{gA}=\mu_{gB}.
\]
可以使用类似Fisher精确检验的方法检验差异表达。考虑表4.6的 $2\times2$ 列联表，其中 $n_g=n_{gA}+n_{gB}$ 为基因 $g$ 的边际行总数，$n=n_A+n_B$ 为总计数。假设从A组抽取 $m_A$ 名受试者，从B组抽取 $m_B$ 名受试者。定义
\begin{equation}
n_{gA}=\sum_{i\in A}y_{gi}=\sum_{i=1}^{m_A}y_{gi}\tag{4.B.42}
\end{equation}
以及
\begin{equation}
n_{gB}=\sum_{j\in B}y_{gj}=\sum_{J=1}^{m_B}y_{gj}.\tag{4.B.43}
\end{equation}
可对表4.6发展类似Fisher精确检验的精确检验，但需以负二项分布概率替代超几何概率。
令 $n_{gA}=a,n_{gB}=b$，以 $p(a,b)$ 表示任意一对数值 $a,b$ 所对应事件的概率。给定总和为 $n_g$ 时，一对实测读段数 $(n_{gA}^*,n_{gB}^*)$ 的 $P$ 值定义为所有不大于实测事件概率 $p(n_{gA}^*,n_{gB}^*)$ 的概率 $p(a,b)$ 之和（Gonzalez 2014）：
\begin{equation}
P_g=\sum_{\substack{a+b=n_g\\p(a,b)\le p(n_{gA}^*,n_{gB}^*)}}p(a,b),\tag{4.B.44}
\end{equation}
其中 $a,b=0,1,\ldots,n_g$。

现在研究如何计算 $p(a,b)$。无表达差异的零假设下，假设不同样本的计数（表达水平）独立，则
\[
p(a,b)=p(n_{gA}=a)p(n_{gB}=b).
\]
可以证明，若 $y_{gi}\sim\mathrm{NB}(r_i,p)$，则
$n_{gA}\sim\mathrm{NB}(\sum_{i=1}^{m_A}r_i,p)$（习题4）。
回顾 $y_{gi}$ 服从负二项分布
$\mathrm{NB}(\phi_{gA}^{-1},\phi_{gA}\mu_{gA}/(1+\phi_{gA}\mu_{gA}))$。
由式（4.B.42）与上述结论，可知 $n_{gA}$ 服从
$\mathrm{NB}(m_A\phi_{gA}^{-1},\phi_{gA}\mu_{gA}/(1+\phi_{gA}\mu_{gA}))$。
类似地，$n_{gB}$ 服从
$\mathrm{NB}(m_B\phi_{gA}^{-1},\phi_{gA}\mu_{gA}/(1+\phi_{gA}\mu_{gA}))$。

利用式（4.B.27），计算得到
\begin{equation}
P(n_{gA}=a)=\frac{\Gamma(a+m_A/\phi_{gA})}{a!\Gamma(m_A/\phi_{gA})}
(\phi_{gA}\mu_{gA})^a
\left(\frac1{1+\phi_{gA}\mu_{gA}}\right)^{a+m_A/\phi_{gA}},\tag{4.B.45}
\end{equation}
以及
\begin{equation}
P(n_{gB}=b)=\frac{\Gamma(b+m_B/\phi_{gB})}{b!\Gamma(m_B/\phi_{gB})}
(\phi_{gB}\mu_{gB})^b
\left(\frac1{1+\phi_{gB}\mu_{gB}}\right)^{b+1/\phi_{gB}}.\tag{4.B.46}
\end{equation}
